\documentclass[11pt]{article}

\usepackage[margin=0.75in]{geometry}
\usepackage{booktabs}
\usepackage{amsmath}
\usepackage{makecell}
\usepackage{hyperref}

\title{Bandwidth-Freeze-Epoch Ablation: freeze=1 vs.\ freeze=5}
\author{}
\date{\today}

\begin{document}
\maketitle

\section{Setup}

Source = syn\_large\_1920 (pretrained, detached, 6400 images, all four weather
conditions), target = real\_small, evaluated on real val (target domain). Hardware,
optimizer, and all non-MMD hyperparameters are identical to those described in
Section~1 of \texttt{experiment\_summary\_v6.tex} (2$\times$RTX 4090, batch 16,
imgsz 1920, 200 max epochs / patience 100, RBF kernel, EMA bandwidth momentum 0.9,
\texttt{mmd\_target\_layer=10}). The only variable changed within each pair below is
the epoch at which the EMA kernel bandwidth stops updating and freezes
(\texttt{bandwidth-freeze-epoch}): 1 (frozen almost immediately) vs.\ 5 (a few
epochs of bandwidth adaptation before freezing). no\_mmd and naive\_mmd rows are
included as unregularized reference points; they do not have a freeze-epoch setting
since naive\_mmd never freezes the bandwidth at all.

\begin{table}[h!]
    \centering
    \small
    \setlength{\tabcolsep}{6pt}
    \begin{tabular}{lrrrr}
        \toprule
        Model & \makecell{person\\mAP50} & \makecell{person\\mAP50-95} & \makecell{vehicle\\mAP50} & \makecell{vehicle\\mAP50-95} \\
        \midrule
        no\_mmd (reference) & 0.648 & 0.296 & 0.870 & 0.617 \\
        naive\_mmd (reference) & 0.637 & 0.291 & 0.869 & 0.618 \\
        \midrule
        reg\_mmd, freeze=1 & 0.635 & 0.292 & 0.871 & 0.619 \\
        reg\_mmd, freeze=5 & 0.642 & 0.294 & 0.872 & 0.618 \\
        \midrule
        reg\_mmd\_smaller\_weight, freeze=1 & 0.639 & 0.291 & 0.871 & 0.617 \\
        reg\_mmd\_smaller\_weight, freeze=5 & 0.635 & 0.289 & 0.876 & 0.620 \\
        \bottomrule
    \end{tabular}
    \caption{syn\_large\_1920 $\to$ real\_small, same regularization schedule and
    weight targets, only bandwidth-freeze-epoch differs within each pair. Overall
    (all-class) mAP is omitted here -- it is a weighted average across both classes
    and is not the quantity of interest for this comparison.}
\end{table}

\section{Discussion}

For reg\_mmd, freeze=5 beats freeze=1 on every column: person mAP50 +0.7
(0.635 $\to$ 0.642), person mAP50-95 +0.2 (0.292 $\to$ 0.294), vehicle mAP50 +0.1
(0.871 $\to$ 0.872), vehicle mAP50-95 roughly flat (0.619 vs.\ 0.618). Freezing the
bandwidth at epoch 1 gives the EMA estimate almost no time to adapt to the actual
feature-space geometry before it locks -- the epoch-5 estimate is evidently a better
fixed point to align towards for the rest of training, and the gain shows up mostly
on person rather than vehicle.

For reg\_mmd\_smaller\_weight the picture is more mixed: freeze=5 gives the best
vehicle score in either ablation table in this study (mAP50 0.876, mAP50-95 0.620)
at a small cost to person (mAP50 0.635 vs.\ 0.639 at freeze=1, mAP50-95 0.289 vs.\
0.291). This mirrors the person/vehicle trade-off already documented in
\texttt{experiment\_summary\_v6.tex} Section~2.2 -- letting the bandwidth adapt a
little longer before freezing appears to push the alignment slightly further toward
the class that dominates the target domain by instance count (vehicle), at a small
cost to the minority class (person).

Overall, freeze=5 is the better default of the two settings tested here: for reg\_mmd
it improves both classes outright, and for reg\_mmd\_smaller\_weight it trades a
small person regression for the best vehicle result in the study. freeze=5 is the
setting used for Section~1 and the newer sections of this study (Sections~2.2
and~3 of \texttt{experiment\_summary\_v6.tex}). freeze=1 remains the setting used
only for the clearnoon side of Section~2.1 in that document, which predates this
ablation; that inconsistency is called out explicitly here rather than implied to
be equivalent.

\end{document}
